% Part: incompleteness
% Chapter: theories-computability
% Section: oconsis-ext-of-q-undec

\documentclass[../../../include/open-logic-section]{subfiles}

\begin{document}

\olfileid{inc}{tcp}{oqn}

\olsection{د~$\Th{Q}$، $\omega$-سازګارې غځونې ناپرېکړه‌وړې دي}

\begin{explain}
دا ثبوت چې~$\Th{Q}$ په محاسبوي ډول د شمېر وړ بشپړه ده، پر دې حقيقت
تکيه درلوده چې په~$\Th{Q}$ کښې هره ثابتېدونکې جمله د طبيعي عددونو
په اړه «رښتيا» ده۔ راتلونکے تعريف او قضيه دا پايله پياوړې کوي، او
کټ مټ د «رښتياوالي» هغه اړخونه ټاکي چې په پورتني ثبوت کښې ورته
اړتيا وه۔ خو پر دې قضيه ډېر مه درېږئ، ځکه ژر به ئې لا نوره هم
پياوړې کړو۔ موږ دا زياتره د تاريخي موخې دپاره راوړو: د ګوډل اصلي
مقالې د~$\omega$-سازګارۍ مفهوم کارولے و، خو لږ وروسته ئې پايله د
$\omega$-سازګارۍ پر ځاے د عادي سازګارۍ په کارولو پياوړې شوه۔
\end{explain}

\begin{defn}
\ollabel{thm:oconsis-q}
يوه تيوري~$\Th{T}$ هله~$\omega$-سازګاره ده چې لاندې شرط پوره کړي:
که~$\lexists[x][!A(x)]$ هره جمله وي او~$\Th{T}$،
$\lnot !A(\num 0)$، $\lnot !A(\num 1)$، $\lnot !A(\num 2)$،
\dots ثابتوي، نو~$\Th{T}$، $\lexists[x][!A(x)]$ نۀ ثابتوي۔
\end{defn}

\begin{thm}
  پرېږدئ~$\Th{T}$ هره هغه~$\omega$-سازګاره تيوري وي چې~$\Th{Q}$
  پکې شامله وي۔ بيا~$\Th{T}$ د پرېکړې وړ نۀ ده۔
\end{thm}

\begin{proof}
که~$\Th{Q}$ په~$\Th{T}$ کښې شامله وي، نو~$\Th{T}$ محاسبه کېدونکې
تابعې او اړيکې تمثيلوي۔ يوازې پخوانی ثبوت لږ بدلول پکار دي۔ لکه
پورته، که~$x \in K$ وي، نو~$\Th{T}$،
$\lexists[s][!A_T(\num x,\num x, s)]$ ثابتوي۔ برعکس، فرض کړئ
$\Th{T}$، $\lexists[s][!A_T(\num x, \num x, s)]$ ثابتوي۔ بيا بايد
$x$ په~$K$ کښې وي: که نۀ وي، د ماشين~$x$ داسې محاسبه نشته چې پر
ننوت~$x$ ودرېږي؛ ځکه~$!A_T$ د کلېني اړيکه~$T$ تمثيلوي، $\Th{T}$،
$\lnot !A_T(\num x, \num x, \num 0)$،
$\lnot !A_T(\num x, \num x, \num 1)$، \dots ثابتوي، او په دې سره
$\Th{T}$، $\omega$-ناسازګاره کېږي۔
\end{proof}

\end{document}
